\label{sec:discussion}

The results support four main conclusions.

\paragraph{Constitutive AD is fastest in the fixed high-order
Plasticity3D derivative test.}
\Cref{fig:plasticity3d-derivative-ablation,tab:plasticity3d-derivative-ablation}
show a clean cost-quality picture on the fixed
$P_4(L_1)$, $\lambda_{\mathrm{sr}}=1.5$, 8-rank case. All three routes
reach essentially the same terminal state with the same nonlinear and Krylov
counts, so the comparison is about cost rather than about different converged
solutions. In this fixed test, constitutive AD preserves the scalar-energy
implementation contract while reducing wall time substantially relative to full
element AD and colored SFD.

\paragraph{Validation has to be separated by claim level.}
The Plasticity3D validation ladder in
\Cref{fig:plasticity3d-validation,tab:plasticity3d-validation} supports a
scoped mechanics interpretation. The tested endpoint surrogate agrees closely
with the source-family observables reported here, and the fixed-$\lambda_{\mathrm{sr}}$
reference-operator comparison satisfies the endpoint field-agreement thresholds
after matching the glued-bottom boundary contract. The evidence therefore
supports the intended endpoint surrogate for the tested cases, but it does not
establish path-consistent incremental elastoplastic-history equivalence or a
verified numeric reproduction of the Sysala source-family papers.

\paragraph{Matched external baselines are useful when the contract is narrow.}
The hyperelastic companion comparison against JAX-FEM in
\Cref{fig:hyperelasticity-jaxfem-baseline,tab:hyperelasticity-jaxfem-baseline}
is informative because the mesh, prescribed displacement schedule, and
terminal-state energy evaluation are explicitly matched. The terminal states
are post-compared with the energy in
\cref{eq:hyperelasticity-energy}; the resulting energy and field differences
are very small. The comparison is useful precisely because it is narrow and
hyperelastic-only. The paper should not generalize it into a broad ranking
between frameworks or into plasticity validation, but it does strengthen the
manuscript by showing that at least one recent external implementation can be
matched on a shared hyperelastic problem contract.

\paragraph{Scope of the evidence.}
The same evidence also defines the scope of the conclusions. The auxiliary
Plasticity3D strong
scaling study at $\lambda_{\mathrm{sr}}=1.0$ remains timing evidence, while
the newer glued-bottom $\lambda_{\mathrm{sr}}=1.55$ scaling sweep is the
converged comparison tied to the benchmark. The topology study is
reported as end-to-end workflow scaling rather than fixed-work strong scaling.
Some reference-operator PMG studies remain solver-engineering diagnostics rather than
fully symmetric scientific baselines. These boundaries give the proper
interpretation of the contribution: a software-and-methods workflow with
explicit derivative-route choices, explicit validation surfaces, and explicit
agreement limits.

The broader methodological lesson is that differentiable FEM becomes most
scientifically useful when local differentiation and sparse solver engineering
are designed together. Automatic differentiation alone does not settle the
behavior of large nonlinear FEM workflows; globalization, sparse ownership,
overlap replication, coarse-level policy, preconditioning, and the interpretation
of surrogate models matter just as much. Within that scoped claim, the toolset
contributes a documented nonlinear FEM workflow rather than a universal
differentiable-PDE platform or a new constitutive theory.
